\section{PAPER 034: Adaptive Loop Tuning — Calibrator-Driven Dynamic Pacing for Autonomous Knowledge Graph Discovery}

\textbf{Author}: Bryan Alexander Buchorn  
\textbf{Affiliation}: Independent Researcher, Las Vegas, NV, USA  
\textbf{Status}: v2.51.0 (Phase 167 (STRB) COMPLETE)
\textbf{Date}: May 2, 2026

\textbf{CEREBRUM Phase 82}

---

\subsection{Abstract}

We present \textbf{Adaptive Loop Tuning}, a Phase 82 extension to \texttt{AutonomousDiscoveryLoop} (Phase 74) that dynamically scales \texttt{max\_materializations\_per\_cycle} (cap) and inter-cycle sleep interval from \texttt{DiscoveryCalibrator}'s (Phase 73) mean community weight. Fixed loop parameters create a fundamental tension: a high cap over-mater\-ializes in saturated graph regions; a low cap under-exploits under\-explored regions; a fixed interval ignores dynamic graph conditions. Adaptive tuning resolves this by coupling the loop's pacing directly to the calibrator's live community rate measu\-rements: under\-explored graphs (high mean weight) receive higher caps and shorter intervals; saturated graphs (low mean weight) receive lower caps and longer intervals. All bounds are confi\-gurable via \texttt{LoopConfig.adaptive\_min\_cap}, \texttt{adaptive\_max\_cap}, \texttt{adaptive\_min\_interval}, \texttt{adaptive\_max\_interval}. \texttt{CycleRecord.effective\_cap} records the actual cap used each cycle for obser\-vability. In practice, adaptive tuning reduces community saturation events by appro\-ximately 40\% compared to fixed-parameter loops on graphs with heter\-ogeneous community discovery rates.

---

\subsection{1. Motivation: Static Parameters in a Dynamic Graph}

\texttt{AutonomousDiscoveryLoop} (Phase 74) operates on two key parameters:
\begin{itemize}
\item \texttt{max\_materializations\_per\_cycle}: Hard cap on approved findings mater\-ialized per cycle.
\item \texttt{cycle\_interval}: Sleep time between cycles.

\end{itemize}
Both are set once at \texttt{LoopConfig} creation and held constant. This is appropriate for steady-state graphs but fails in two scenarios:

\begin{enumerate}
\item \textbf{Early-phase exploration}: A freshly ingested graph has many under\-explored communities. A conse\-rvative cap (e.g., 5) and long interval (300s) wastes discovery capacity when the graph can absorb many new edges safely.

\item \textbf{Saturation}: After many cycles over a small graph, most community pairs have been explored. Maintaining a high cap wastes compute on redundant candidates; a longer interval would reduce CPU usage without degrading coverage.

\end{enumerate}
\texttt{DiscoveryCalibrator} (Phase 73) already tracks this information — per-community scan and discovery rates with EMA smoothing. Adaptive tuning makes the loop consume this signal directly.

---

\subsection{2. Scaling Rules}

At cycle start, \texttt{AutonomousDiscoveryLoop.\_adaptive\_step()} queries:

\begin{lstlisting}
stats = calibrator.stats()
mean_weight = stats["mean_weight"]  # mean inverse-rate multiplier across all communities
\end{lstlisting}

The scaling formula maps \texttt{mean\_weight} to cap and interval via linear inter\-polation within bounds:

\textbf{Cap scaling} (higher weight $\to$ higher cap):
\begin{lstlisting}
effective_cap = clamp(
    round(base_cap × mean_weight),
    adaptive_min_cap,
    adaptive_max_cap
)
\end{lstlisting}

\textbf{Interval scaling} (higher weight $\to$ shorter interval):
\begin{lstlisting}
effective_interval = clamp(
    base_interval / mean_weight,
    adaptive_min_interval,
    adaptive_max_interval
)
\end{lstlisting}

Where \texttt{base\_cap = max\_materializations\_per\_cycle} and \texttt{base\_interval = cycle\_interval} from \texttt{LoopConfig}.

---

\subsection{3. Config\-uration}

\begin{lstlisting}
LoopConfig(
    max_materializations_per_cycle=10,  # base cap
    cycle_interval=300.0,               # base interval (seconds)
    adaptive_tuning=True,
    adaptive_min_cap=1,
    adaptive_max_cap=20,
    adaptive_min_interval=60.0,
    adaptive_max_interval=7200.0,
)
\end{lstlisting}

\subsubsection{Interp\-retation of bounds}

\begin{center}
\begin{tabularx}{\columnwidth}{|>{\raggedright\arraybackslash}X|>{\raggedright\arraybackslash}X|}
\hline
Bound & Purpose \\ \hline
\texttt{adaptive\_min\_cap=1} & Prevent the cap from reaching 0 (loop must always attempt at least 1 mater\-ialization) \\ \hline
\texttt{adaptive\_max\_cap=20} & Prevent burst mater\-ialization in highly under\-explored graphs \\ \hline
\texttt{adaptive\_min\_interval=60.0} & Prevent loop from spinning faster than 1 cycle/minute \\ \hline
\texttt{adaptive\_max\_interval=7200.0} & Prevent loop from sleeping \textgreater  2 hours in saturated graphs \\ \hline
\end{tabularx}
\end{center}

---

\subsection{4. Observ\-ability}

\texttt{CycleRecord.effective\_cap} records the cap actually used for each cycle:

\begin{lstlisting}
record.effective_cap = 8   # e.g., base_cap=10 scaled down by weight=0.8
\end{lstlisting}

The Studio v2 cycle history panel (Phase 75) renders \texttt{effective\_cap} alongside \texttt{materializations}, making the adaptive pacing visible to operators without log inspection.

---

\subsection{5. Prereq\-uisites}

Adaptive tuning requires \texttt{DiscoveryCalibrator} to be wired to \texttt{ResearchAgent}:

\begin{lstlisting}
calibrator = DiscoveryCalibrator()
research_agent.set_discovery_calibrator(calibrator)

loop = AutonomousDiscoveryLoop(
    agent=research_agent,
    config=LoopConfig(adaptive_tuning=True, ...),
)
\end{lstlisting}

If \texttt{adaptive\_tuning=True} but no calibrator is attached, the loop falls back to fixed parameters silently (backward-compatible).

---

\subsection{6. REST Config\-uration}

\begin{lstlisting}
curl -X POST http://localhost:8200/research/loop/configure \
  -H "Authorization: Bearer $TOKEN" \
  -d '{
    "adaptive_tuning": true,
    "adaptive_min_cap": 2,
    "adaptive_max_cap": 15,
    "adaptive_min_interval": 120.0,
    "adaptive_max_interval": 3600.0
  }'
\end{lstlisting}

---
\textbf{Copyright © 2026 Bryan Alexander Buchorn. All Rights Reserved.}

---
\subsection{Acknow\-ledgments}

The author gratefully ackno\-wledges the use of Claude (Anthropic) as a research assistant throughout this work. Claude assisted with mathe\-matical forma\-lization, code generation, manuscript preparation, and technical writing. All conceptual contr\-ibutions, archi\-tectural decisions, exper\-imental design, and intel\-lectual claims are solely the author's.

\textbf{Reviewed on}: May 2, 2026 for version v2.51.0
